\documentclass[11pt]{article}
\usepackage[margin=1in]{geometry}
\usepackage{amsmath,amssymb,amsthm}
\usepackage{booktabs}
\usepackage{hyperref}
\usepackage{xcolor}

\newtheorem{proposition}{Proposition}
\newtheorem{lemma}[proposition]{Lemma}
\theoremstyle{remark}
\newtheorem{remark}[proposition]{Remark}

\newcommand{\R}{\mathbb{R}}
\newcommand{\Sym}{\mathcal{S}}
\newcommand{\tr}{\operatorname{tr}}
\newcommand{\vecop}{\operatorname{vec}}
\newcommand{\Diag}{\operatorname{Diag}}
\newcommand{\inner}[2]{\langle #1, #2\rangle}
\newcommand{\lmax}{\lambda_{\max}}
\newcommand{\Pspd}{\mathcal{P}_{+}}
\newcommand{\QPB}{\mathrm{QPB}}
\newcommand{\DNN}{\mathrm{DNN}}
\newcommand{\GL}{\mathrm{GL}}
\newcommand{\LAP}{\mathrm{LAP}}
% filled in from the clean nug30 run (experiments/nug30_solve)
\newcommand{\CleanNodes}{6,532,419}
\newcommand{\CleanCPU}{18,614}
\newcommand{\CleanWall}{$\approx$1.3 h}
\newcommand{\CleanDNN}{the DNN bound ran on 5,764 nodes, fathomed 5,551, and accounted for 17,192 CPU seconds}

\title{A certified doubly nonnegative bound for branch and bound\\on the quadratic assignment problem,\\and nug30 revisited}
\author{Nathan W. Brixius\thanks{Draft technical report, October 2026. Source and data:
the \texttt{cognac} branch of \url{https://github.com/natebrix/qap}. See the
acknowledgments for how this report and its software were produced.}}
\date{October 2026}

\begin{document}
\maketitle

\begin{abstract}
In 2000, the quadratic assignment problem nug30 was solved by a branch-and-bound
algorithm based on a convex quadratic programming bound (QPB), running for about
a week on a computational grid of hundreds of workstations. We revisit nug30 with
the same QPB-based branch and bound, strengthened near the top of the tree by a
much stronger, much more expensive bound: the doubly nonnegative (DNN)
semidefinite relaxation of Zhao, Karisch, Rendl and Wolkowicz, in the facially
reduced form of Oliveira, Wolkowicz and Xu, solved by ADMM. Because a
branch-and-bound proof may only use valid bounds, and an ADMM iterate is not
one, we compute at every check a \emph{certified} lower bound: the Lagrangian
dual value of a multiplier whose projection onto the relaxation's face is
negative semidefinite by construction, evaluated by a Gilmore--Lawler-type
minimization over permutations, less a floating-point safety margin built from
standard error bounds (for the spectral part, Rump's Cholesky criterion). We
describe the bound, its certificate, an $O(m^4)$-per-product implementation of
the facially reduced ADMM with a single-precision subspace eigensolver, and its
use in branch and bound, where it is applied only to nodes at depths 2--6 that QPB
does not fathom, to fathom them or raise their bounds; branching and reduced-cost
elimination keep using QPB. With it, nug24 is solved in 299 nodes and 13 CPU minutes, and
nug30 in 6.5 million nodes, 5.2 CPU hours and about 80 minutes of wall time on
four cores.
\end{abstract}

\section{Introduction}

The quadratic assignment problem (QAP) in Koopmans--Beckmann form is
\begin{equation}\label{eq:qap}
  \min_{X \in \Pi_n} \; \tr(A^T X B X^T) + \tr(C X^T)
  = \min_{\sigma} \sum_{i,k} a_{ik} b_{\sigma(i)\sigma(k)} + \sum_i c_{i\sigma(i)},
\end{equation}
where $\Pi_n$ is the set of $n\times n$ permutation matrices ($X_{i\sigma(i)}=1$),
$A$ (flows) and $B$ (distances) are $n\times n$, and $C$ is a linear cost. (When $A$
is symmetric, as for the Nugent instances, the first term is the more familiar
$\tr(AXBX^T)$.) QAPLIB \cite{qaplib} has long
served as its benchmark library, and the Nugent instances \cite{nugent} nug12 to nug30
as its yardstick for exact methods. nug30 resisted solution for three decades until
Anstreicher, Brixius, Goux and Linderoth \cite{abgl} solved it in 2000 by branch and
bound on the quadratic programming bound QPB of Anstreicher and Brixius \cite{ab01},
on a computational grid of several hundred workstations over about a week (see
\cite{abgl} for the exact node count and CPU time).

QPB is a bound of moderate strength that is cheap to compute approximately by
Frank--Wolfe iterations, each requiring one linear assignment problem (LAP). That
balance suited a search of billions of nodes. Semidefinite relaxations are far
stronger: the DNN relaxation \cite{zkrw,povhrendl} is within about 3\% of the optimum at
the root of nug30, against about 13\% for QPB (Table~\ref{tab:root}). But they have
been considered too expensive to apply at more than a few nodes. Fujii et al.\
\cite{fujii} showed that a Lagrangian-DNN branch and bound, with every node bounded
by the DNN relaxation, can solve nug30-sized instances with tens of thousands of
nodes and several previously unsolved QAPLIB instances, on large parallel
computers. Graham et al.\ \cite{graham} derive the same spectral dual bound we use
(Proposition~\ref{prop:cert}) for a strengthened DNN relaxation solved by a
restricted Peaceman--Rachford splitting.

This report describes the opposite division of labor: QPB does the bulk of the
search, and a certified DNN bound is applied only where QPB fails at the top of the
tree, at depths 2--6. Nodes there are few, but each carries a subtree of
$10^3$--$10^9$ QPB nodes. On nug30 the DNN bound costs seconds per node and fathoms
almost all of them. QPB's value at a node predicts well whether the DNN bound will
fathom it.

The contributions are:
\begin{itemize}
\item a certificate (Section~\ref{sec:cert}) that turns any ADMM multiplier into a
  lower bound, valid in exact arithmetic and with explicit floating-point error
  allowances under stated assumptions, cheap enough to evaluate every few
  iterations, so that the bound can stop as soon as it fathoms: the multiplier
  implied by the $R$-step, and a Gilmore--Lawler evaluation of the polyhedral part;
\item an implementation (Section~\ref{sec:impl}) of the facially reduced ADMM with
  $O(m^4)$ products by the face basis and a single-precision tracked eigensolver,
  13--36$\times$ faster per node than a numpy prototype;
\item its integration into a QPB branch and bound (Section~\ref{sec:bnb}), and
  computational results on nug20--nug30 (Section~\ref{sec:results}), including
  nug30 solved in 5.2 CPU hours, about 80 minutes on four cores.
\end{itemize}

\section{The DNN relaxation}\label{sec:dnn}

Consider a node of the branch-and-bound tree with $m$ facilities left to assign.
Its subproblem is again a QAP \eqref{eq:qap} of size $m$, with reduced matrices
$A, B \in \R^{m\times m}$, a linear term $C$ that collects the interactions with
the fixed assignments, and a constant $s$ (the solver's \texttt{reduceQAP}). For
$X\in\Pi_m$ let $x = \vecop(X)$, ordered so that $x_{ij}$ (facility $i$ at location
$j$) has lifted index $1 + im + j$, and let $y = (1, x)$ and $N = m^2+1$.
Define the symmetric lifted cost $L\in\Sym^N$ by
\begin{equation}\label{eq:L}
  L_{(ij),(kl)} = \tfrac12(a_{ik}b_{jl} + a_{ki}b_{lj}), \qquad
  L_{0,(ij)} = L_{(ij),0} = \tfrac12 c_{ij}, \qquad L_{00} = 0.
\end{equation}
Then for every permutation, $\inner{L}{yy^T} = \sum a_{ik}b_{jl}x_{ij}x_{kl} +
\sum c_{ij}x_{ij} = \tr(A^TXBX^T) + \tr(CX^T)$, the subproblem objective without
$s$, for any $A$ and $B$, symmetric or not.

\paragraph{The face.} The constraints $Xe = X^Te = e$ say that $y$ lies in the null
space of a linear map, so every lifted permutation $Y = yy^T$ has its range in a
fixed subspace of dimension $(m-1)^2+1$. Let $V \in\R^{N\times K}$, $K = (m-1)^2+1$,
be an orthonormal basis of it. Zhao et al.\ \cite{zkrw} take
$V = \begin{pmatrix}1 & 0\\ e/m & V_m\otimes V_m\end{pmatrix}$, suitably normalized,
where the columns of $V_m\in\R^{m\times(m-1)}$ span $e^\perp$. Every lifted permutation then
satisfies $Y = VRV^T$ with $R \succeq 0$. The DNN relaxation of the subproblem is
\begin{equation}\label{eq:dnn}
  \DNN = s + \min\{ \inner{L}{Y} : Y = VRV^T,\ R\succeq 0,\ Y_{00}=1,\
  \mathcal{G}(Y) = 0,\ 0 \le Y \le 1 \},
\end{equation}
where $\mathcal{G}$ selects the ``gangster'' entries $Y_{(ij),(il)}$, $l\ne j$, and
$Y_{(ij),(kj)}$, $k\ne i$, which vanish for permutations since a facility has one
location and a location one facility. Oliveira, Wolkowicz and Xu \cite{owx} solve
\eqref{eq:dnn} by ADMM on the splitting $Y = VRV^T$, alternating
\begin{align}
  R &\leftarrow \Pspd\big(V^T(Y + Z/\beta)V\big), \label{eq:rstep}\\
  Y &\leftarrow \Pi_{\mathcal{Y}}\big(VRV^T - (L + Z)/\beta\big), \label{eq:ystep}\\
  Z &\leftarrow Z + \gamma\beta\,(Y - VRV^T), \label{eq:zstep}
\end{align}
where $\Pspd$ projects onto the positive semidefinite cone,
$\mathcal{Y} = \{Y: Y_{00}=1, \mathcal{G}(Y)=0, 0\le Y\le1\}$ is a box with fixed
entries (so $\Pi_{\mathcal{Y}}$ is entrywise clipping), $\beta>0$ is the penalty
and $\gamma\in(0,\frac{1+\sqrt5}{2})$ the dual step.

\section{A certified bound from any multiplier}\label{sec:cert}

ADMM iterates are neither primal feasible nor dual optimal, and the primal value
$\inner{L}{Y}$ is not a bound. Branch and bound needs a number that is provably
below the subproblem optimum after any number of iterations. We obtain one from
the multiplier.

\begin{lemma}\label{lem:trace}
Every lifted permutation $Y = yy^T$ satisfies $Y = VRV^T$ with
$R = V^TYV\succeq0$ and $\tr(R) = m+1$.
\end{lemma}
\begin{proof}
$y$ lies in the range of $V$, so $Y = VV^TYVV^T = VRV^T$ with $R = V^TYV =
(V^Ty)(V^Ty)^T\succeq0$, and $\tr R = \|V^Ty\|^2 = \|y\|^2 = 1 + \sum x_{ij}^2 =
m+1$.
\end{proof}

\begin{proposition}\label{prop:cert}
Let $Z\in\Sym^N$ be any symmetric matrix, $\lambda \ge \lmax(V^TZV)$, and let
$\mathcal{P}\supseteq\{yy^T : X\in\Pi_m\}$. Then the subproblem optimum is at least
\begin{equation}\label{eq:cert}
  s + \min_{Y\in\mathcal{P}} \inner{L+Z}{Y} - (m+1)\lambda .
\end{equation}
\end{proposition}
\begin{proof}
For a lifted permutation $Y = VRV^T$ (Lemma~\ref{lem:trace}),
$\inner{Z}{Y} = \inner{V^TZV}{R} \le \lmax(V^TZV)\tr R \le (m+1)\lambda$, since $R\succeq0$.
Hence $\inner{L}{Y} = \inner{L+Z}{Y} - \inner{Z}{Y} \ge \inner{L+Z}{Y} - (m+1)\lambda
\ge \min_{\mathcal{P}}\inner{L+Z}{\cdot} - (m+1)\lambda$.
\end{proof}

The same bound appears as the dual function of Graham et al.\ \cite[eq.~(3.11),
Thm.~3.4]{graham}. Proposition~\ref{prop:cert} holds for every $Z$; the quality of the bound depends
on $Z$ being near a dual optimum of \eqref{eq:dnn}, and on choosing $Z$ so that
$\lambda$ can be small. Two choices of $\mathcal{P}$ and one of $Z$ make it both
cheap and tight.

\paragraph{The multiplier.} The prototype of \cite{owx} projects $V^TZV$ onto the
negative semidefinite cone, which costs an extra eigendecomposition per check.
Instead we use the multiplier implied by the $R$-step. Let $W = V^T(Y+Z/\beta)V$ be
the matrix projected in \eqref{eq:rstep}, and
\[
  Z' = Z + \beta(Y - VRV^T), \qquad\text{so that}\qquad V^TZ'V = \beta(W - \Pspd(W)) \preceq 0 .
\]
The last matrix is $\beta$ times the negative part of $W$, negative semidefinite by
construction. This holds even when the eigenvectors used for $\Pspd$ are
inaccurate: if $\Pspd(W)$ is formed as $W$ minus the computed negative part
$\sum_{w_k<0} w_k v_kv_k^T$, the difference is exactly that sum. In floating
point, $\lmax(V^TZ'V)$ is a tiny positive number, which we bound rigorously below.

\paragraph{The polyhedral part.} With $\mathcal{P}$ the box
$\mathcal{Y}$ itself, $\min_{\mathcal{P}}\inner{M}{Y}$ for $M = L+Z'$ is
$M_{00}$ plus the sum of the negative entries of $M$ outside the gangster set: the
prototype's choice. A much tighter, still cheap choice uses that every permutation
is an assignment in two levels. Write $M_{(ij),\cdot}$ for row $(ij)$ of $M$. For a
permutation $\sigma$,
\begin{equation}\label{eq:gl}
\inner{M}{yy^T} = M_{00} + \sum_{i}\Big( 2M_{0,(i\sigma_i)} + M_{(i\sigma_i),(i\sigma_i)}
 + \sum_{k\ne i} M_{(i\sigma_i),(k\sigma_k)} \Big).
\end{equation}
The inner sum is a perfect matching between facilities $k\ne i$ and locations
$l\ne\sigma_i$, so it is at least the value $\LAP_{ij}$ of the $(m-1)\times(m-1)$
assignment problem on the block $M_{(ij),(kl)}$, $k\ne i$, $l\ne j$. Hence
\begin{equation}\label{eq:glbound}
\min_{\sigma}\inner{M}{yy^T} \;\ge\; \GL(M) := M_{00} + \LAP\big(c\big),\qquad
c_{ij} = 2M_{0,(ij)} + M_{(ij),(ij)} + \LAP_{ij},
\end{equation}
a Gilmore--Lawler computation \cite{gilmore,lawler} on the lifted matrix. (It also
bounds the DNN relaxation itself, not only permutations: for feasible $Y$ of
\eqref{eq:dnn}, column $(ij)$ restricted to $k\ne i$, $l\ne j$ is nonnegative with
all row and column sums $x_{ij} = Y_{0,(ij)}$, so it is $x_{ij}$ times a doubly
stochastic matrix, and assignment duality applies twice.) It is
never below the box value and reaches a given value 1.5--6$\times$ sooner in the
nug20 tests. Any lower bound on each $\LAP_{ij}$ may replace it, for example the
row-then-column reduction, which costs $O(m^2)$. The implementation uses the
reduction, plus exact Hungarian solves when the result is within a small gap of the
cutoff.

\paragraph{Floating point.} We assume IEEE binary64 arithmetic with rounding to
nearest, an extended type (\texttt{long double}) with a $p$-bit significand
($p=64$ on x86), and no overflow or underflow; a non-finite value anywhere yields
no certificate. The certified value is $s + \GL(L+Z') - \mu$, rounded down, where
the inputs and the margin $\mu$ cover:
\begin{itemize}
\item[(o)] the subproblem data: $s$ and the reduced linear term $C$, sums of
  products that can cancel, are replaced by lower bounds (exact when all operands
  are multiples of $\frac12$ and the sum of magnitudes is below $2^{p-2}$, as for
  integer data; otherwise the extended-precision sum less $(k+2)\,2^{-p}\sum|t|$
  for $k$ terms $t$, rounded down). Bounds for the lowered subproblem are bounds for
  the exact one, since assignment variables are nonnegative. Each entry of $L$
  is charged $u$ times the sum of its two products' magnitudes.
\end{itemize}
\begin{itemize}
\item[(i)] $(m+1)\lambda$ with $\lambda\ge\lmax(V^TZ'V)$, obtained by Rump's
  criterion \cite{rump}: if the floating-point Cholesky factorization of
  $\delta I - S$ succeeds for the computed $S \approx V^TZ'V$, then
  $\lmax(S) \le \delta + c$ with an explicit $c$ depending on the dimension, the
  unit roundoff $u$ and $\tr(\delta I - S)$. To this is added
  $64\,m\,u\,\|Z'\|_F$ for the error of computing $S$ through Householder
  transforms rather than with an exact $V$.
\item[(ii)] the summation and rounding errors of \eqref{eq:glbound}, accumulated in
  extended precision: $8(m+2)u\,\|M\|_1$, plus $(m+2)u\sum_{(ij),(kl)}
  \frac12(|a_{ik}b_{jl}|+|a_{ki}b_{lj}|)$ for forming and scaling $L$.
\end{itemize}
These are standard error bounds with generous constants, not interval arithmetic,
and not every constant is derived here: the allowance $64\,m\,u\|Z'\|_F$ for the
Householder transforms rests on a standard backward-error argument and was checked
empirically (errors at most $2.1\,u\|Z'\|_F$), not proved.
At $m=28$ the largest margin seen was about $10^{-7}$ in objective units, against
integer objective values. The solver fathoms only when the certified value exceeds
the incumbent minus $1.95$ for instances whose objective values are provably even
(Section~\ref{sec:bnb}).

\begin{remark}
Nothing in the certificate depends on the ADMM having converged, on the accuracy
of its eigenvectors, or on the single-precision arithmetic used in the $R$-step
(Section~\ref{sec:impl}). Only the multiplier $Z'$, the Cholesky test and
\eqref{eq:glbound} enter the certified value, all in double precision.
\end{remark}

\section{Implementation}\label{sec:impl}

The bound is implemented in C++ (\texttt{latest/dnn.C}, about 1100 lines) with
LAPACK and BLAS (OpenBLAS). A numpy prototype (\texttt{experiments/dnn\_bound}) was
ported faithfully; a prototype mode reproduces its certified values to $10^{-7}$.

\paragraph{Products with $V$.} The prototype forms $V$ explicitly and computes
$V^TMV$ and $VRV^T$ with dense products of order $N\times N$, which costs
$O(m^6)$. We take for $V_m$ the first $m-1$ columns of the Householder
reflector $H$ with $He_{m-1} = e/\sqrt m$. Then $V^TMV$ and $VRV^T$ are a
reflection along each of the four tensor modes plus a $2\times2$ rotation coupling
the constant coordinate, $O(m^4)$ per product. Both the $R$-step and the
certificate are invariant under an orthogonal change of basis of the face, so
any orthonormal $V$ gives the same iterates and bound.

\paragraph{Eigensolver.} Each $R$-step needs the projection of a
$K\times K$ symmetric matrix ($K = 730$ at $m=28$) with hundreds of eigenvalues of
each sign, so a partial spectrum does not help directly. Between iterations the
matrix changes little. We track its negative eigenspace by a Chebyshev-filtered
subspace iteration in single precision (BLAS level 3), with a full single-precision
eigendecomposition every 20 iterations. Because $R$ is formed as $W$ minus the
computed negative part, the certificate is unaffected (Section~\ref{sec:cert}).
When a check fails the Cholesky test, or the semidefinite penalty alone decides
whether the node is fathomed, the multiplier is re-projected in double precision
before certifying.

\paragraph{Parameters.} The objective is scaled to $\max|L| = 1$. We use
$\beta = 0.3\cdot m/3$ (the prototype's $m/3$ times 0.3), which certifies
2--3$\times$ sooner on hard nodes, and $\gamma = 1.618$. A certificate is computed
every 10 iterations, and the run stops at the first certified value above the
cutoff. An optional give-up rule stops a node early when its certified value,
extrapolated generously from its recent gains, stays below the cutoff. This only
affects cost: the node is then branched on as usual.

\paragraph{Speed.} Table~\ref{tab:speed} compares the C++ bound with the prototype
on the same nodes. Measured back to back on a loaded 4-core Xeon, an ADMM iteration
at $m=28$ costs 43\,ms against 339\,ms. Per node, the C++ bound needs fewer
iterations to certify, thanks to the tighter polyhedral part and the smaller $\beta$,
and is 13--36$\times$ faster.

\begin{table}[t]\centering
\caption{The certified DNN bound in C++ against the numpy prototype, CPU seconds per
node, cutoff = optimum $-1.95$ (one core; prototype times from the studies in
\texttt{experiments/dnn\_bound}).}\label{tab:speed}
\begin{tabular}{lrrrr}
\toprule
node set & fathomed (both) & C++ & prototype & ratio\\
\midrule
nug20, depth 2 (all 98, $m=18$) & 96 & 0.62 & 8.35 & 13$\times$\\
nug30, depth 2 (36 random, $m=28$) & 34 & 9.6 & 139 & 14.5$\times$\\
nug30, depth 2 (8 lowest QPB) & 0 & 11 (200 it.) & 397 & 36$\times$\\
nug30, depth 3 (20, $m=27$) & 19 & 7.4 & 121 & 16$\times$\\
\bottomrule
\end{tabular}
\end{table}

\section{Use in branch and bound}\label{sec:bnb}

The branch and bound is the QPB code of \cite{ab01,abgl}, with improvements made
since then: Frank--Wolfe warm-started from the parent's solution, keeping the best
(bound, reduced cost) pair over all iterations, and cheaper, corrected lookahead
branching at the top levels. For grid instances, it also uses a branching rule on
the most eccentric free location down to depth 6 (\texttt{PARAM/nug\_deep.par}).
The DNN bound is added as follows.

\begin{enumerate}
\item At a node of depth $d\in[d_1,d_2]$ whose QPB value does not fathom it, the
  certified DNN bound is computed for its reduced subproblem
  (\texttt{bnbdrv -x d1:d2[:gap[:giveup]]}).
\item If it exceeds the incumbent minus the fathoming tolerance, the node is
  fathomed. Otherwise it raises the node's recorded bound, which its children
  inherit and which also feeds the choice of parameters and some branching
  scores, so the search differs from a QPB-only search even where the DNN bound
  does not fathom.
\item Branching, and the elimination of children by $\text{bound} + U_{ij} \ge
  \text{incumbent}$, use the QPB value and its reduced costs $U$, never the DNN
  value. The elimination argument is a property of the Frank--Wolfe linearization
  of QPB.
\end{enumerate}

The fathoming tolerance is 1.95 when all objective values are provably even, as for
the nug instances, 0.95 when they are integers, and 0 otherwise; data is treated
as integer only if every objective value is below $2^{53}$, so that it is computed
exactly. For even values and an even incumbent, an improving solution is at most
the incumbent minus 2, so 1.95 leaves a cushion of 0.05 for rounding.

\paragraph{Scope.} The DNN certificate covers only the DNN prunes. QPB values, QPB
reduced-cost eliminations and symmetry reductions are computed in floating point
without such margins, as in \cite{abgl}, so the search as a whole is not an
independently checkable proof. The full solver also needs one of $A$, $B$ to be
symmetric (otherwise QPB is not valid, and the driver warns), and its QPB search is
currently wrong on data with negative entries in $A$ or $B$ (Section~\ref{sec:validation}).
The DNN bound has neither restriction. Search statistics,
including DNN counts, are kept in the checkpoint files, so a run can be stopped,
resumed and split across processes (\texttt{bench/parallel.py}) with identical
totals. One caveat: the DNN bound's floating-point results depend on the BLAS
kernels selected for the CPU, so node counts with the DNN bound can differ between
machines if a certified value lands within about $10^{-6}$ of a cutoff. Validity is
not affected.

\paragraph{Where to apply it.} Two studies with the prototype measured the bound
on QPB-open nodes (Table~\ref{tab:fathom}). At depth 1 it rarely fathoms. At depth 2
it fathoms nearly all nodes of nug20, and 34 of a random sample of 36 of nug30's 225
(94\%). In that sample the outcome was almost a function of the node's QPB value:
every node with QPB below about 5548 was missed, every node above about 5554
fathomed. Applied to the whole depth-2 frontier, that threshold predicts about
185 of 225 fathoms (82\%), a sample-based estimate rather than a guarantee for
other depths or instances. At
depth 3, under the depth-2 nodes it misses, it fathoms almost everything. The
misses are the most expensive subtrees, so the bound has to be applied at several
consecutive depths. In the full solves, depths 2 to 5 or 6 worked best.

\begin{table}[t]\centering
\caption{Root bounds and DNN fathoming at the top of the tree (prototype studies;
QPB from the solver; incumbent = optimum).}\label{tab:root}\label{tab:fathom}
\begin{tabular}{lrrr}
\toprule
& nug12 & nug20 & nug30\\
\midrule
optimum & 578 & 2570 & 6124\\
root QPB & 478.5 & 2243.9 & 5332.3\\
root DNN (certified) & 568.0 & 2506.3 & 5949.3\\
\midrule
depth 1: fathomed / QPB-open & -- & 1 / 6 & 0 / 1 sampled\\
depth 2: fathomed / QPB-open & -- & 96 / 98 & 34 / 36 sampled\\
depth 3: fathomed / QPB-open & -- & 200 / 200 sampled & 19 / 20 sampled\\
\bottomrule
\end{tabular}
\end{table}

\section{Computational results}\label{sec:results}

All runs use the known optimum as the incumbent, as in \cite{abgl}, so they measure
the work of the proof. They ran on a cloud virtual machine with 4 cores (Intel Xeon,
2.1--2.8\,GHz), with \texttt{bench/parallel.py} distributing subtrees to 4 solver
processes. CPU is the sum over processes.

\begin{table}[t]\centering
\caption{Branch and bound with and without the DNN bound (4 workers unless noted).
QPB only: default parameters. DNN: \texttt{-x d1:d2}, giving up on hopeless nodes
after 300 iterations where shown (:300).}\label{tab:bnb}
\begin{tabular}{llrrr}
\toprule
instance & bounds & nodes & CPU s & wall\\
\midrule
nug20 & QPB only (1 core) & 526,205 & 197 & --\\
nug20 & QPB + DNN 2--3 (1 core) & 141 & 105 & --\\
nug24 & QPB only & $\approx$25\,M & $\approx$5,000 & $\approx$25 min\\
nug24 & QPB + DNN 2--3 & 299 & 796 & 4.4 min\\
nug25 & QPB only & 52,989,998 & 18,353 & 78 min\\
nug25 & QPB only (\texttt{nug\_deep}) & 40,183,792 & 12,362 & 53 min\\
nug25 & QPB (\texttt{nug\_deep}) + DNN 2--5 & 717 & 3,279 & 15 min\\
nug25 & QPB (\texttt{nug\_deep}) + DNN 2--5:300 & 761 & 1,447 & 6.2 min\\
nug27 & QPB (\texttt{nug\_deep}) + DNN 2--5 & 675 & 4,428 & 20 min\\
nug28 & QPB (\texttt{nug\_deep}) + DNN 2--5 & 1,736 & 12,976 & 62 min\\
nug30 & QPB (\texttt{nug\_deep}) + DNN 2--5:300, then 2--6:300 & 45,486,332 & 40,555 & $\approx$3.0 h\\
nug30 & QPB (\texttt{nug\_deep}) + DNN 2--6:300 & \CleanNodes & \CleanCPU & \CleanWall\\
\midrule
\multicolumn{5}{l}{\emph{on a second machine, about 1.45$\times$ slower per DNN iteration:}}\\
nug30 & QPB (\texttt{nug\_deep}) + DNN 2--6:300 & 6,532,419 & 28,428 & $\approx$2.1 h\\
nug30 & QPB (\texttt{nug\_deep}) + DNN 2--7:300, warm & 5,817 & 23,618 & 1.7 h\\
\bottomrule
\end{tabular}
\end{table}

Table~\ref{tab:bnb} summarizes the solves. At nug20 the DNN bound at depths 2--3
removes all but 141 of 526,205 nodes and halves the time. At nug24 the tree
collapses to 299 nodes. From nug25 on, the DNN bound does most of the work: on
nug28, 12,926 of 12,976 CPU seconds are spent in it, on 1,727 nodes of which 1,664
are fathomed.

The two nug25 rows with \texttt{nug\_deep} and the 2--5 DNN setting isolate the
DNN bound: the same parameters, branching rules and symmetry handling, with the
bound on or off. Other comparisons in the table change parameters as well.

\paragraph{nug30.} The first solve applied the DNN bound at depths 2--5. QPB search
below the depth-5 nodes the bound missed grew to about 12,000 CPU seconds, so for
the remaining work (506 open nodes, including a whole depth-1 subtree) the bound
was extended to depth 6, which finished in 21 minutes. In total, 45,486,332 nodes
and 40,555 CPU seconds. The DNN bound ran on 5,615 nodes, fathomed 5,403, and
accounted for 26,770 of those seconds. A clean run with the DNN bound at depths 2--6
throughout took \CleanNodes{} nodes, \CleanCPU{} CPU seconds and \CleanWall{} of wall
time (\CleanDNN): seven times fewer nodes and less than half the CPU of the
first solve. Almost all of the difference is QPB search below depth-6 nodes that
the depth-5 setting left to QPB. Nearly all (92\%) of the clean run's time is in
the DNN bound, and 96\% of the nodes it bounds are fathomed. The run started on a
2.1\,GHz Xeon and, after a container restart, finished its last 10 minutes on a
2.8\,GHz one. The CPU figures are retained proof CPU: work done after the last
checkpoint before a restart (6 minutes of the clean run's 4 workers; 8 minutes in
the first solve) was discarded and redone, and is not counted. The clean run used
solver commit \texttt{f8107dd} with OpenBLAS (\texttt{scipy-openblas32}) and GCC.
Later fixes from the reviews below change no computation on nug data, whose
products and sums are exact; the regression instances' node counts are
unchanged.

\paragraph{Improved settings.} Two changes target the clean run's remaining costs.
The QPB search below its 72 depth-7 nodes took about 1,300 CPU seconds; applying the
DNN bound at depth 7 too fathoms all of them (96 in that run, 178 seconds). And most
DNN calls are at children of nodes whose own DNN bound missed; starting a child's
ADMM from the parent's final multiplier, mapped to the child's lifted indices,
cuts the iterations to fathom by about 40\%. The run with both
(\texttt{-x 2:7:1:300:1}) ran on a different machine, so the clean settings were
rerun there with the same code (the last two rows of Table~\ref{tab:bnb}). The
rerun reproduces the clean run's tree exactly, and the new settings take 17\% less
CPU, with 5,817 nodes in all: every node at depths 2--7 that QPB does not fathom
gets the DNN bound, and all of them are fathomed by depth 7. Warm starts lower the
DNN time per node at depths 3 and 4 by only 16\% and 13\%, less than the iteration
counts suggest, because a warm-started ADMM pays for many negative eigenvalues
from its first iteration where a cold one starts cheap (on nug25 the saving was
4\%).

\section{Validation}\label{sec:validation}

A branch-and-bound proof is only as good as its bounds, so the DNN bound was
checked in several independent ways.
\begin{itemize}
\item \textbf{Exhaustive checks.} 1,512 certified bounds on random instances and
  subproblems with $m\le8$ (several families, including linear terms and ties),
  each compared with the optimum found by enumeration. No bound exceeds the
  optimum, and no fathom is ever claimed with the cutoff at the optimum.
\item \textbf{Inside the search.} The solver's brute-force test (4,536 runs on
  random instances with $n\le 8$, eight parameter sets, three upper-bound modes)
  passes with the DNN bound applied at every depth, with and without the give-up
  rule.
\item \textbf{Against the prototype.} In prototype mode the C++ bound reproduces the
  numpy prototype's certified values at 44 checkpoints to a relative $8.6\cdot10^{-8}$.
  On nug20 nodes whose subproblems were solved exactly, the bounds are below the
  optima and fathom the same nodes.
\item \textbf{Reproducibility.} Results are bitwise identical across runs on one
  machine. Valgrind reports no memory errors on a debug build. Serial, resumed and
  parallel runs give identical node counts.
\item \textbf{Independent mathematical review.} A separate AI agent, given the code
  and this claim but not our arguments, re-derived every step: the face basis
  actually used (Householder reflector and rotation) has the required range, so
  Lemma~\ref{lem:trace} holds for it; Proposition~\ref{prop:cert} holds for any
  sign of $\lmax$; the Gilmore--Lawler evaluation counts every quadratic term
  once and is valid for any dual potentials, so neither the inexact inner LAPs
  nor the outer Hungarian solve need to be exact; the lifting matches the
  solver's objective convention without symmetry assumptions; a NaN or infinite
  multiplier cannot produce a fathom; and the search uses the DNN value only to
  fathom. It found no proof-breaking issue.
\item \textbf{Independent code review and stress tests.} A second agent reviewed
  the implementation and ran 2,592 bounds on brute-forced instances with
  $m\le7$ (asymmetric and negative data, ties, scales up to $10^6$), 540
  warm-started children, an independent long-double check of the face
  transforms (error at most $2.1\,u\|Z\|_F$ against the $64\,m\,u\|Z\|_F$
  allowed), 21,000 random assignment problems against the Hungarian solver,
  and determinism tests, under address, undefined-behavior and leak sanitizers.
  No bound exceeded an optimum and nothing crashed. Its minor findings (and the
  mathematical review's) were fixed: a loop that could skip the last Cholesky
  attempt, widened error constants, an explicit finiteness test on the
  certified value, and skipping the bound while there is no incumbent.
\item \textbf{Adversarial tests.} A third agent tried to break the bound with
  extreme data: magnitudes up to $10^{9}$ and beyond, products above $2^{53}$,
  signed, non-integer and non-symmetric data, rank-one, constant and all-zero
  matrices, subproblems with one to three free facilities, and parameter
  fuzzing; about 15,000 bounds and 8,600 solver runs in all, with an
  instrumented build checking every DNN call inside the search against the
  exact subproblem optimum. On ordinary data no bound was invalid, and the
  rounding part of the margin was never needed (the unmargined value minus the
  semidefinite penalty stayed below the optimum by at least 1.6\% of the rounding
  margin). It did find two ways to exceed the optimum, both needing data that is
  neither symmetric nor small. When neither $A$ nor $B$ is symmetric, an entry of
  $L$ can be the difference of two large products whose rounding errors the
  margin, which scaled with $\sum|L|$, did not cover; the margin now uses the sum
  of the products' magnitudes. And the shift and linear term of a subproblem, sums
  of products that can cancel, were rounded in double before any bound saw them.
  Integer data whose products exceed $2^{53}$ is no longer given an integer
  fathoming tolerance. It also found
  two older bugs outside the DNN bound: a finite sentinel ($10^8$) in the
  assignment solver, which crashed or gave invalid QPB bounds when costs exceeded
  it (fixed), and wrong QPB results on data with negative entries in $A$ or $B$
  (open; QAPLIB data are nonnegative). None of this affects nonnegative integer
  data whose objective values stay below $2^{53}$, where every product and sum is
  exact. That includes all of QAPLIB, even the bur26 family, in which neither
  matrix is symmetric. On the regression instances, node counts are unchanged.
\item \textbf{External review.} An independent review of an earlier draft and the
  code (Astra, an OpenAI Codex agent) re-derived the bound and confirmed that QPB
  values stay paired with their own reduced costs. It showed that accumulating the
  shift and linear terms in extended precision, our first repair, is not enough:
  with $ab = 2^{64}+2^{34}+3$ and $c = -(2^{64}+2^{34})$ the extended-precision
  product rounds by 1 before cancelling, and both the $m=1$ path and a reduced
  subproblem fathomed at a cutoff equal to the optimum. These data now get the
  lower bounds of item (o) in Section~\ref{sec:cert}; both cases, and the
  adversarial ones above, are regression tests (\texttt{tests/dnn\_test.py}, part D)
  that check the bound against optima computed in exact rational arithmetic.
\end{itemize}

\paragraph{Limitations.} The certificate is a theorem in exact arithmetic
(Proposition~\ref{prop:cert} and \eqref{eq:glbound}). In floating point it rests
on the margins of Section~\ref{sec:cert}: standard error bounds with generous
constants under the stated arithmetic model, not interval arithmetic, and not
every constant is derived. The validation above is strong regression evidence, not
a formal verification. For the nug instances the data are small integers, so the
subproblem data are exact, and a false fathom would need a certified value above
the optimum by more than the 0.05 cushion, five to six orders of magnitude above
any margin observed. For non-integer data, where the tolerance is zero, validity
rests on the margins alone; a verified evaluation (directed rounding or interval
arithmetic) of each fathoming certificate would be the safe choice. Beyond the DNN
bound, the QPB part of the search is trusted floating point, as it was in 2000.

\section{Discussion}

Twenty-six years ago nug30 needed QPB's cheapness: billions of nodes, each bounded
in milliseconds, on hundreds of machines. Today the same search needs only a few
thousand much stronger bounds at the top, and QPB below them, to finish on one small
machine in an afternoon. Hardware accounts for part of the gain, but most of it
comes from the division of labor. The stronger bound, certified at every step, is
spent only where a single fathom removes a large subtree.

Several directions remain. The give-up rule and the depths of application were
tuned by hand. A QPB threshold predicts the DNN outcome almost perfectly in the
samples, which suggests skipping nodes the bound will miss. One idea from the
review did not pay off on nug30: the outer assignment problem of \eqref{eq:glbound}
yields, with its dual potentials $u,v$, a lower bound $b + r_{ij}$ for every child
$x_{ij}=1$ of a node the bound misses, with $r_{ij} = c_{ij}-u_i-v_j$, at no extra
cost (\texttt{dnntest -C}). On nine sampled depth-1 and depth-2 misses it fathomed
none of their 784--841 children each: near a dual optimum the outer costs are almost flat (median
$r_{ij}$ 0.1, largest 12.1, where the nearest miss needed 14.4), and even forcing
the assignment in the outer problem gains at most 12.1. A child's gain comes from
re-solving the relaxation with the extra fixing. The restricted
Peaceman--Rachford method of \cite{graham} may converge faster than ADMM, and a
separate verified check of each fathoming certificate would separate fast search
from rigorous verification.

\paragraph{Open instances.} We screened candidates among the instances QAPLIB lists
as open by their certified root DNN gap (nug30: 2.9\%, and tho30, solved in
\cite{abgl}: 4.2\%): wil50 1.4\%, sko49 3.1\%, tho40 5.8\%, tai50b 6.3\%, tai35a
8.5\%, tai40a 9.4\%. At depth 2, 7 of 10 random nug30 nodes are fathomed, but
none of 10 for tho40 and none of the sko49 nodes sampled, whose gap is still about
2\%. At $n\approx50$ an ADMM iteration costs about 1 second on four threads,
13$\times$ nug30's, and the bound would have to reach depth 3--4 before fathoming,
which means $10^4$--$10^6$ DNN nodes. The open grid instances are out of reach
on one small machine with this method; they would need much larger parallel
resources, as in \cite{fujii}, or a faster DNN solver.

\section*{Acknowledgments}
The software described here, the experiments and a first draft of this report were
produced with Claude, an AI system by Anthropic, working in the author's repository
under the author's direction, with independent AI agents used for the prototype
studies, the C++ port and the reviews. Astra (an OpenAI Codex agent) reviewed an
earlier draft and found the cancellation cases described in
Section~\ref{sec:validation}. Jean-Pierre Goux contributed the ``nug30
Laboratory'' study whose ideas started this work.

\begin{thebibliography}{99}
\bibitem{ab01} K.M. Anstreicher and N.W. Brixius, A new bound for the quadratic
  assignment problem based on convex quadratic programming, \emph{Mathematical
  Programming} 89 (2001) 341--357.
\bibitem{abgl} K. Anstreicher, N. Brixius, J.-P. Goux and J. Linderoth, Solving large
  quadratic assignment problems on computational grids, \emph{Mathematical
  Programming} 91 (2002) 563--588.
\bibitem{qaplib} R.E. Burkard, S.E. Karisch and F. Rendl, QAPLIB -- a quadratic
  assignment problem library, \emph{Journal of Global Optimization} 10 (1997) 391--403.
\bibitem{fujii} K. Fujii, N. Ito, S. Kim, M. Kojima, Y. Shinano and K.-C. Toh, Solving
  challenging large scale QAPs, arXiv:2101.09629 (2021).
\bibitem{gilmore} P.C. Gilmore, Optimal and suboptimal algorithms for the quadratic
  assignment problem, \emph{SIAM Journal on Applied Mathematics} 10 (1962) 305--313.
\bibitem{graham} N. Graham, H. Hu, J. Im, X. Li and H. Wolkowicz, A restricted dual
  Peaceman--Rachford splitting method for a strengthened DNN relaxation for QAP,
  \emph{INFORMS Journal on Computing} 34 (2022) 2125--2143.
\bibitem{lawler} E.L. Lawler, The quadratic assignment problem, \emph{Management
  Science} 9 (1963) 586--599.
\bibitem{nugent} C.E. Nugent, T.E. Vollmann and J. Ruml, An experimental comparison of
  techniques for the assignment of facilities to locations, \emph{Operations
  Research} 16 (1968) 150--173.
\bibitem{owx} D.E. Oliveira, H. Wolkowicz and Y. Xu, ADMM for the SDP relaxation of
  the QAP, \emph{Mathematical Programming Computation} 10 (2018) 631--658.
\bibitem{povhrendl} J. Povh and F. Rendl, Copositive and semidefinite relaxations of
  the quadratic assignment problem, \emph{Discrete Optimization} 6 (2009) 231--241.
\bibitem{rump} S.M. Rump, Verification of positive definiteness, \emph{BIT
  Numerical Mathematics} 46 (2006) 433--452.
\bibitem{zkrw} Q. Zhao, S.E. Karisch, F. Rendl and H. Wolkowicz, Semidefinite
  programming relaxations for the quadratic assignment problem, \emph{Journal of
  Combinatorial Optimization} 2 (1998) 71--109.
\end{thebibliography}

\end{document}
